\section{Capítulo~\ref{sec:eetypes}. Neurônios como componentes}

Os modos de funcionamento de células individuais foram medidos diretamente, com corrente por eletrodo e registro de tensão; a matemática do integrador e da divisão em classes é teoria clássica, e o uso da reconfiguração de modos pelo cérebro para calcular continua sendo uma hipótese.

{\footnotesize
\setlength{\tabcolsep}{3pt}\renewcommand{\arraystretch}{1.15}
\begin{longtable}{>{\raggedright}p{0.5cm}>{\raggedright}p{4.0cm}>{\raggedright}p{5.6cm}>{\raggedright}p{2.3cm}>{\raggedright\arraybackslash}p{1.7cm}}
\toprule
N.º & Proposição & Experimento e o que foi medido & Fonte & Status \\
\midrule\endfirsthead
\toprule
N.º & Proposição & Experimento e o que foi medido & Fonte & Status \\
\midrule\endhead
1 & Ressonador: uma corrente lenta que se opõe à mudança de tensão torna o neurônio um filtro passa-faixa com máximo entre alguns hertz e dezenas de hertz (seção~\ref{sec:resonator}) & Em fatias, injetou-se em neurônios \nt{GLUT} corrente senoidal com frequência crescente e mediu-se a impedância $|Z(f)|$: máximo em $2$--$5$~Hz a $33^\circ$C (cerca de $7$~Hz a $38^\circ$C); a ressonância é criada por correntes lentas de potássio e catiônica e amplificada pela corrente persistente de sódio. Uma revisão mostra o mesmo recurso em muitas classes de células & \cite{HuVervaekeStorm2002,HutcheonYarom2000} & medido \\
2 & A corrente lenta se comporta como indutância e, com a capacitância da membrana, forma um circuito ressonante & A resposta sublimiar do axônio à corrente foi medida e descrita por equações de condutância linearizadas; o circuito equivalente contém uma indutância ``fenomenológica'' cujo valor depende da tensão & \cite{Mauro1970} & teoria \\
3 & Constante de tempo do grupo com realimentação $\tau_{\text{ef}}=\tau/(1-w)$~\eqref{eq:perfectint}; com $\tau=0.1$~s e $w=0.995$, $20$~s & Deduzida no capítulo a partir da equação linear do grupo; proposta como mecanismo de manutenção da posição do olho num trabalho teórico & dedução no capítulo; \cite{Seung1996} & teoria \\
4 & O integrador da posição do olho mantém a entrada pela rede, e não por uma propriedade de uma célula isolada & Registro intracelular em peixe, no núcleo que mantém a posição do olho: após um giro rápido, a taxa do neurônio muda em degrau e se mantém; uma corrente breve numa célula causa só um deslocamento passageiro, e os degraus de tensão persistem abaixo do limiar: são mantidos pelas entradas sinápticas & \cite{Aksay2001} & medido \\
5 & O integrador sem vazamento precisa de ajuste contínuo por aprendizado; desajustado, ele esquece ou satura & Peixes treinados com realimentação visual proporcional a $\pm$ a posição do olho: a fixação ficava instável ou com vazamento, e a constante de tempo dos neurônios caía mais de uma ordem de grandeza, para $<1$~s; a realimentação visual normal restabelecia aos poucos a estabilidade & \cite{Major2004} & medido \\
6 & Neurônio biestável: uma entrada curta liga um platô duradouro, a inibição o desliga; a propriedade é ligada pela serotonina (seção~\ref{sec:bistable}) & Registro intracelular de neurônios \nt{ACET} que controlam músculos, em gato: uma excitação sináptica curta leva o neurônio a funcionamento prolongado, uma inibição curta o reinicia; no gato espinhal, o estado aparece após a administração de um precursor da serotonina & \cite{Hounsgaard1988} & medido \\
7 & Duas classes: integradores (a taxa cresce a partir do zero) e ressonadores (no limiar, frequência finita de imediato) & As classes decorrem da teoria de bifurcações. Experimento: em fatias de córtex, neurônios \nt{GLUT} começam a disparar com taxa tão baixa quanto se queira (tipo~1), e os neurônios \nt{GABA} rápidos, com salto para uma frequência finita (tipo~2) & \cite{Izhikevich2007,Tateno2004} & medido \\
8 & Um neurônio é integrador ou detector de coincidência: a inibição encurta a janela de soma & Em fatias de hipocampo, a janela em que as entradas excitatórias se somam até o limiar é menor que $2$~ms; ela é encurtada pela inibição direta que chega logo após a excitação; nos dendritos, onde a inibição é mais fraca, a janela é mais larga & \cite{Pouille2001} & medido \\
9 & Linha de atraso feita de axônios (tabela~\ref{tab:eetypes}) & Na coruja-das-torres foram medidos os atrasos nos axônios que vão aos detectores de coincidência: comprimentos diferentes de axônio dão atrasos diferentes, e os detectores estão sintonizados com a diferença dos tempos de chegada do som & \cite{Carr1990} & medido \\
10 & Marca-passo na vigília e célula silenciosa no sono & Os neurônios \nt{SERO} funcionam de modo quase regular de dia, desaceleram no sono de ondas lentas e quase se calam no sono REM (mais detalhes no capítulo~\ref{sec:sleep}) & \cite{JacobsAzmitia1992,McGintyHarper1976} & medido \\
11 & O componente é definido pelos canais iônicos e pelos canais de difusão que os ajustam (proposição~\ref{prop:eetypes}) & Mostrado em redes pequenas: substâncias moduladoras alteram as propriedades intrínsecas dos neurônios e a força das sinapses, de modo que o mesmo circuito produz ritmos diferentes & \cite{Marder2012} & medido \\
12 & O cérebro usa a reconfiguração de modos para calcular (proposição~\ref{prop:eetypes}) & Não há experimento direto em que a reconfiguração do modo de uma célula seja necessária para resolver uma tarefa; há só experimentos em que os moduladores mudam a saída do circuito & \cite{Marder2012} & hipótese \\
\bottomrule
\end{longtable}}
